\documentclass[11pt]{article}

\usepackage[margin=1in]{geometry}
\usepackage{amsmath}
\usepackage{booktabs}
\usepackage{float}
\usepackage{graphicx}
\usepackage{hyperref}
\usepackage{natbib}
\usepackage{url}
\Urlmuskip=0mu plus 1mu\relax

\title{Projected Fixed Points and Internal Recurrent Dynamics:\\A Measurement Protocol in Demian}
\author{Azael\\Independent Researcher}
\date{June 13, 2026}

\begin{document}

\maketitle

\begin{abstract}
Demian is an experimental discrete-time nonlinear recurrent system with full
state $z_t$ and exposed readout $y_t$. We test the hypothesis that $y_t$ can
converge while $z_t$ retains structured, continuation-relevant dynamics.
Dual-GRU experiments show identical fixed-point surface labels with different
internal regimes, including tight and accumulating classes. We define an
accumulating fixed point operationally by surface convergence together with
persistent structured hidden-state change, after excluding numerical drift,
ordinary transients, and trivial accumulation. Channel ablations and matched
perturbation baselines bound the interpretation. Full-state versus
surface-only continuation is used only as a sufficiency control: full restore
contains more information, and the observed gap establishes that the exposed
state is not sufficient for continuation. Jacobian and finite-time Lyapunov
analysis remain future characterization.
\end{abstract}

\noindent\textbf{Keywords:} recurrent neural networks; dynamical systems;
fixed points; hidden state; mechanistic interpretability; state-space models.

\section{Introduction}

Demian is the experimental system, not the hypothesis. We write its dynamics
as
\[
  z_{t+1}=F(z_t,x_t), \qquad y_t=R(z_t),
\]
where $z_t$ is the full recurrent state and $y_t$ is the exposed readout.

The core hypothesis is that $y_t$ may converge while $z_t$ retains structured,
continuation-relevant dynamics. A fixed point of the readout is therefore not
necessarily a fixed point of the full system.

The paper tests that hypothesis in dual-GRU and native Demian artifacts. Plain
recurrent, Transformer, and Mamba runs provide controls or architectural
context, not a benchmark ranking. The target is whether one surface class
contains systematically different internal regimes and whether those regimes
align with explicit recurrent channels.

The contribution is a compact evidence discipline:

\begin{itemize}
  \item classify the exposed surface trajectory;
  \item separately classify fixed-point interiors;
  \item measure route and channel activity;
  \item perturb or ablate internal state;
  \item compare full internal-state resume against surface-only resume;
  \item demote mechanism names when strict held-out checks fail.
\end{itemize}

\begin{figure}[H]
\centering
\includegraphics[width=0.92\linewidth]{figures/surface_internal_protocol.png}
\caption{Protocol-level distinction used throughout the paper. The exposed
surface can be classified first, but the fixed-point label is not treated as a
verdict until internal-state probes, ablations, and resume controls are
checked.}
\label{fig:protocol}
\end{figure}

\section{Background and related work}

The fixed-point vocabulary used here comes from the dynamical-systems view of
recurrent computation. In this view, the system state evolves under a repeated
map and fixed points, cycles, and transient regimes are useful summaries of
that evolution. Such summaries are valuable because they compress complex
trajectories into interpretable regimes, but they can also hide differences
between systems that expose similar readouts.

Prior work has used fixed points and line-attractor structure to
reverse-engineer trained recurrent networks \citep{maheswaranathan2019lineattractor}.
Other work shows that networks can solve the same task with different internal
solutions \citep{huang2024degeneracy}, and that comparing dynamics requires
tools that are sensitive to trajectory geometry rather than only static
representations \citep{guilhot2024dsa}. Modern sequence architectures are also
increasingly analyzed as dynamical systems, including Transformer dynamics
\citep{fernando2025transformerdynamics} and selective state-space models such
as Mamba \citep{gu2023mamba}. This paper fits inside that line of work by
asking a smaller operational question: when a recurrent surface appears fixed,
what measurements are needed before concluding that the internal computation
has collapsed?

\section{Terminology}

The full recurrent state is $z_t$ and the exposed surface is $y_t=R(z_t)$. A
surface fixed point satisfies $\Delta y_t \to 0$ under the experiment's
convergence tolerance. It does not imply $\Delta z_t \to 0$.

An \emph{accumulating fixed point} (AFP) is assigned only when surface
convergence coexists with persistent structured change in $z_t$. The analysis
must distinguish that pattern from floating-point drift, a finite transient,
and trivial unstructured accumulation.

A \emph{channel} is a named component of $z_t$. A \emph{surface-only
continuation} initializes hidden state to zero or omission while restoring
$y_t$. A \emph{full-state continuation} restores $z_t$. Their comparison is a
sufficiency control, not a mechanism discovery test.

We use five epistemic labels: \emph{observation} for artifact-backed
measurements, \emph{hypothesis} for falsifiable relations, \emph{control} for
comparisons that bound sufficiency, \emph{interpretation} for bounded
inference, and \emph{untested speculation} for mechanisms not yet measured.


\section{Methods}

\subsection{Substrates and baselines}

The analysis follows this order:

\begin{enumerate}
  \item classify surface convergence from $y_t$;
  \item measure persistence and change in the full state $z_t$;
  \item compare internal regimes within the same surface class;
  \item run full-state versus surface-only continuation as a sufficiency control;
  \item restore or ablate channels and components;
  \item measure perturbation response and stability;
  \item reserve Jacobian and finite-time Lyapunov analysis for the next rigorous characterization.
\end{enumerate}

The artifact source is the Demian Lab repository snapshot tagged
\texttt{v0.2-fixed-point-paper} \citep{demianlab2026}. The paper's tracked
evidence summary is included beside the paper source. The original source
artifacts remain in the repository's \texttt{data/} tree.

The experiments were not designed as a single benchmark suite. They are a
curated set of diagnostic artifacts from the same research line. The seed
counts are therefore reported with each result rather than hidden behind a
single aggregate. Dual-GRU interior maps use eight runs per listed variant.
Matched perturbation baselines use three runs per architecture. The capsule
continuity table reports the summary extrema from the v9 continuity sweep.

\subsection{Surface and interior measurements}

The first measurement layer labels the exposed trajectory as a fixed point,
cycle, transient, or related regime. The second layer asks whether fixed-point
surfaces are internally homogeneous. Dual-GRU summaries report interior
classes, message-state norms, bottleneck code counts, and bottleneck entropy.
These quantities are not presented as universal metrics; they are the local
measurement interface for this substrate family.

\subsection{Continuation-state sufficiency control}

Capsule continuity asks whether a paused trajectory can be resumed from
different reconstructions of state. Let $s_t$ be the internal recurrent state
and $y_t = R(s_t)$ be the exposed surface. The uninterrupted continuation is:
\[
  s_{t+k} = F^k(s_t).
\]
The full capsule control restores $s_t$ and compares its continuation with the
uninterrupted trajectory. The surface-only control constructs a state
$\hat{s}_t$ from $y_t$ alone, leaving the non-surface internal state empty or
zeroed, and compares:
\[
  F^k(\hat{s}_t) \quad \textrm{against} \quad F^k(s_t).
\]
If full-state continuation matches and surface-only continuation diverges, the
exposed surface was not sufficient for continuation. The full restore contains
more information, so this control is expected to favor it and is not treated as
the main discovery.

\subsection{Claim discipline}

The paper separates measurement claims from mechanism names. The gate-state
truth campaign is included because it demoted a stronger mechanism label under
stricter held-out checks. That negative result is part of the method: a
surface/internal difference is evidence for hidden structure, but not
automatically evidence for a named causal mechanism.

\section{Results}

\subsection{Surface convergence and internal-state persistence}

All listed dual-GRU variants converge to the same exposed fixed-point class.
Their hidden message norms and bottleneck statistics remain differentiated.

\subsection{Same surface class, different internal regimes}

Table~\ref{tab:dual-gru} shows the central internal-regime result. Every
listed dual-GRU variant is classified as \texttt{FIXED\_POINT} in 8/8 runs at
the exposed surface. The internal readings are not identical. Some variants
split into both tight and accumulating fixed-point interiors, while
\texttt{dual\_gru\_v3b:current} is entirely accumulating and has much larger
mean message norm and bottleneck entropy.

\begin{table}[H]
\centering
\small
\begin{tabular}{llllrr}
\toprule
Substrate & Surface & Interior counts & Classes & Message norm & Entropy \\
\midrule
\texttt{dual\_gru\_v3b:current} & 8/8 fixed & accumulating: 8 & 1 & 23.1086 & 1.7779 \\
\texttt{dual\_gru\_v3b:tight} & 8/8 fixed & tight: 6; accumulating: 2 & 2 & 0.0624 & 0.3423 \\
\texttt{dual\_gru\_v3b:threshold} & 8/8 fixed & tight: 5; accumulating: 3 & 2 & 0.0753 & 0.5120 \\
\texttt{dual\_gru\_v3} & 8/8 fixed & tight: 6; accumulating: 2 & 2 & 0.0624 & 0.5069 \\
\bottomrule
\end{tabular}
\caption{Dual-GRU fixed-point surfaces with different basin interiors.
All rows share the same exposed surface attractor class, but differ in hidden
interior class, message-state expression, and bottleneck entropy.}
\label{tab:dual-gru}
\end{table}

This is the simplest evidence that a fixed-point label is not enough. The
surface class is constant across the table, while the internal class and
internal activity vary.

\begin{figure}[H]
\centering
\includegraphics[width=\linewidth]{figures/dual_gru_interiors.png}
\caption{Dual-GRU fixed-point interiors. Each variant has the same exposed
surface class, \texttt{FIXED\_POINT} in 8/8 runs, but the measured interior
class count, message norm, and bottleneck entropy differ.}
\label{fig:dual-gru}
\end{figure}

\subsection{Continuation-state sufficiency control}

Table~\ref{tab:capsule} shows capsule-continuity probes for canonical
\texttt{demian\_native\_v9} and the v9 five-channel substrate. Across the
reported sweep, full internal-state resume remains exact or near-exact, while
surface-only resume remains worse than full capsule resume.

\begin{table}[H]
\centering
\small
\begin{tabular}{lrrr}
\toprule
Substrate & Full capsule min cosine & Full capsule max gap & Surface-only min gap \\
\midrule
\texttt{demian\_native\_v9} & 0.99999988 & 0.0 & 0.22624873 \\
\texttt{v9\_five\_channel} & 0.99999988 & 0.0 & 0.27455845 \\
\bottomrule
\end{tabular}
\caption{Capsule-continuity summary. Full internal-state capsules resume
the deterministic continuation, while surface-only reconstruction does not.}
\label{tab:capsule}
\end{table}

This establishes only that the visible surface is not a sufficient
continuation state in these probes. Full restore contains more information than
the surface-only zeroed reconstruction; the comparison does not by itself
establish an AFP or a novel internal mechanism.

\begin{figure}[H]
\centering
\includegraphics[width=0.76\linewidth]{figures/capsule_resume_gap.png}
\caption{Capsule-continuity control. Full internal-state capsules resume the
deterministic continuation with zero measured mean gap in these summaries,
while surface-only reconstruction leaves a continuation gap.}
\label{fig:capsule-gap}
\end{figure}

\subsection{Channel/component restore and ablation}

Component-only restore and channel ablation test whether internal regimes align
with explicit recurrent channels. Current artifacts suggest concentration in
\texttt{slow} for canonical v9 and distribution across \texttt{carrier},
\texttt{message}, and \texttt{slow} in the five-channel system. This is a
channel-alignment hypothesis pending replicated null-controlled tests.

\subsection{Perturbation and stability analysis}

Table~\ref{tab:baselines} summarizes matched pseudo-channel perturbation
baselines from the truth-campaign comparison. Plain RNN, GRU, and LSTM
baselines recover almost exactly under this protocol. Diagonal SSM and native
Demian substrates show larger perturbation consequences.

\begin{table}[H]
\centering
\small
\begin{tabular}{lrrr}
\toprule
Substrate & Runs & Final L2 gap mean & Recovery-window gap mean \\
\midrule
\texttt{rnn} & 3 & $1.19\times10^{-7}$ & 0.000111 \\
\texttt{gru} & 3 & $5.18\times10^{-8}$ & 0.000144 \\
\texttt{lstm} & 3 & $4.46\times10^{-8}$ & 0.000061 \\
\texttt{diag\_ssm} & 3 & 0.0650 & 0.002366 \\
\texttt{demian\_native\_v8} & 3 & 0.2206 & 0.003906 \\
\texttt{demian\_native\_v9} & 3 & 0.00322 & 0.000531 \\
\bottomrule
\end{tabular}
\caption{Matched perturbation/recovery baselines. The table is not a
superiority ranking. It shows why baseline controls are necessary before
promoting internal-structure claims.}
\label{tab:baselines}
\end{table}

The conservative reading is that some perturbation effects are ordinary
recurrent recovery behavior, while other effects require more careful
mechanism tests. This is why the paper centers measurement discipline rather
than global architecture claims.

\subsection{Architecture context: Transformer and Mamba reservoirs}

The Transformer and Mamba reservoir artifacts are useful context but not the
main proof. The Transformer self-reference reservoir supported a deterministic
period-2 attractor observation in the Demian claims ledger. Mamba
self-reference runs did not simply reproduce that period-2 behavior; in the
batch summary, cached Mamba runs have lag-2 autocorrelation 0.1495, while
no-cache runs show 570 period-2 switches in 1000 steps. This supports a
modest point: recurrence-like surface behavior depends on architecture, so
surface classes should not be treated as architecture-independent evidence.

\subsection{Negative gate-state checks keep the claim bounded}

The 2026-05-16 truth campaign demoted the stronger gate-state mechanism name
under strict held-out checks: the strict Track B profile pass rate was 0/2 in
the inspected multisurgery summary. This does not weaken the fixed-point
paper's main claim. It strengthens the paper's boundary. Hidden structure is
not the same as settled causal mechanism naming.

\section{Discussion}

The results support a practical distinction: fixed point is a surface class,
not a computational verdict. In the dual-GRU family, the same surface label
contains different interior classes. In the capsule-continuity probes, a
surface-only reconstruction fails to preserve continuation even when a full
internal capsule resumes exactly. In the matched baseline comparison, plain
recurrent models and native substrates behave differently under the same
pseudo-channel perturbation protocol.

The broader lesson is methodological. A recurrent system should not be judged
only by the final or exposed trajectory. Surface labels should be paired with
state-sensitive controls: interior class, route norms, channel separation,
ablation response, and resume fidelity.

This is also the design reason Demian v1 makes internal channels explicit. The
claim is not that v1 is complete. The claim is that the experimental record
made anonymous hidden state too weak an object of study. If continuation lives
inside the state, then the state must be named, saved, ablated, and tested.

\subsection{Next rigorous characterization}

The current diagnostics do not identify the local stability spectrum of the
full state dynamics. Jacobian analysis and finite-time Lyapunov estimates are
the next tests; until then, chaos, invariant manifolds, and distinct internal
attractors remain untested speculation.

\section{Limitations}

The evidence is heterogeneous because it comes from a living research
repository rather than a single polished benchmark suite. Some runs have small
seed counts. Some metrics are local to Demian's substrate families. The
Transformer and Mamba reservoir artifacts are architectural context rather
than central evidence. The continuation control uses full internal-state restore, not a compressed state code, and therefore contains strictly more information than the surface-only arm. Finally, the paper does not prove that
every fixed point is meaningful; it proves that fixed-point surface labels can
be incomplete.

\section{Reproducibility}

The source repository is Demian Lab \citep{demianlab2026}. The primary
artifact files used in this draft are:

\begin{itemize}
  \item \path{data/substrate_lab/dual_gru_family_summary.json}
  \item \path{data/substrate_lab/v9_capsule_continuity_20260511/summary.json}
  \item gate-state truth campaign multisurgery baseline-comparison summary
  \item gate-state truth campaign multisurgery campaign summary
  \item \path{data/mamba_batch/mamba_batch_summary.json}
  \item \path{data/reservoir_batch/batch_summary.json}
\end{itemize}

The paper-local evidence table is stored beside this source as
\path{evidence_summary.csv} and \path{evidence_summary.json}. The paper can be
built with:

\begin{verbatim}
cd papers/fixed_point_internal_structure
make
\end{verbatim}

The intended public reproducibility anchor is the GitHub tag
\texttt{v0.2-fixed-point-paper}. That tag should point at the final committed
paper source, figure files, and evidence summaries.

\section{Conclusion}

Projected fixed-point behavior is a statement about $y_t$, not necessarily
about $z_t$. The Demian artifacts show the same convergent surface class
coexisting with measurably different internal regimes. The main research target
is now to test whether those regimes align reproducibly with explicit channels
and to characterize their stability with Jacobian and finite-time Lyapunov
analysis. Full-state continuation remains a sufficiency control.

\appendix

\section{Artifact index}

\begin{sloppypar}
\small
The exact artifact references used by the draft are:

\begin{itemize}
  \item \path{papers/fixed_point_internal_structure/evidence_summary.csv}
  \item \path{papers/fixed_point_internal_structure/evidence_summary.json}
  \item \path{data/substrate_lab/dual_gru_family_summary.json}
  \item \path{data/substrate_lab/v9_capsule_continuity_20260511/summary.json}
  \item \path{data/diagnostics/gate_state_truth_campaign_20260516_3trackb_multisurgery/baseline_comparison_summary.json}
  \item \path{data/diagnostics/gate_state_truth_campaign_20260516_3trackb_multisurgery/truth_campaign_summary.json}
  \item \path{data/mamba_batch/mamba_batch_summary.json}
  \item \path{data/reservoir_batch/batch_summary.json}
\end{itemize}
\end{sloppypar}

\bibliographystyle{plainnat}
\bibliography{references}

\end{document}
